\section{The constructive diagonal contradiction}
\label{sec:diagonal}

Combine the analytic witness equivalence of Section~\ref{sec:reduction}
with \eqref{eq:negative-mrdp}. Let
\[
 P_e(\boldsymbol y)=P_0(e,\boldsymbol y),\qquad C(e)=\ulcorner E_{P_e}\urcorner.
\]
Specialization of a fixed polynomial, computation of its coefficient
bounds, and construction of the expression tree are primitive recursive.
The internal calibration theorem is therefore
\begin{equation}\label{eq:calibration}
 \forall e\;\bigl(\Id(C(e))\longleftrightarrow\neg\Halt(e)\bigr).
\end{equation}
Every $C(e)$ is a valid expression in the smaller language, with a single
outer absolute-value occurrence.

\subsection{The exact decider specification}

Write $\operatorname{Bit}(b)$ for $b=0\lor b=1$. Define
\begin{align*}
 \operatorname{Total}(d)&:\Longleftrightarrow
  \forall a\bigl(\WF(a)\to
      \exists s\exists b\,(
          \operatorname{Bit}(b)\land\Run(d,a,s,b))\bigr),\\
 \operatorname{Correct}(d)&:\Longleftrightarrow
  \forall a\forall s\;\Bigl(\WF(a)\to
       \bigl((\Run(d,a,s,1)\to\Id(a))\land
             (\Run(d,a,s,0)\to\neg\Id(a))\bigr)\Bigr),\\
 \Dec(d)&:\Longleftrightarrow
       \operatorname{Total}(d)\land\operatorname{Correct}(d).
\end{align*}
Output $1$ means ``identically zero.'' The correctness condition for
output $0$ is only $\neg\Id(a)$; we do not require the decider to return a
positive nonzero witness. Thus the impossible task is already the weaker
Boolean decision task.

\subsection{The diagonal program and its finite trace laws}

Given a proposed decider code $d$, construct a program $\delta_d$ with the
following behavior on input $z$:
\begin{lstlisting}
a := C(z)
run program d on input a
if the returned output is 1 then halt
otherwise enter a permanent loop
\end{lstlisting}
If the inner program does not terminate, neither does $\delta_d$. The
compiler $d\mapsto\ulcorner\delta_d\urcorner$ is primitive recursive:
it inserts the numeral $d$ into a fixed simulation-and-branching template.
Set $e=\ulcorner\delta_d\urcorner$.

Because $\Halt(e)$ is halting on input $e$, the relevant simulated input
is exactly $C(e)$. Finite trace verification proves
\begin{align}
 \Run(d,C(e),s,1)&\longrightarrow\Halt(e),
       \label{eq:diag-forward}\\
 \Halt(e)&\longrightarrow\exists s\;\Run(d,C(e),s,1).
       \label{eq:diag-backward}
\end{align}
For \eqref{eq:diag-forward}, append the compiler's evaluation trace and
the halt branch to the supplied inner trace. For
\eqref{eq:diag-backward}, inspect a supplied terminating trace of
$\delta_d$ and extract its completed inner simulation. The only halt
instruction reachable in the wrapper is the output-$1$ branch. Both
transformations inspect or generate bounded finite data, and their
correctness uses quantifier-free induction on that data.

No recursion theorem is required. The program accepts a variable input
$z$; after its code is constructed, we evaluate it on that same code.
This is the ordinary diagonal-input convention from
\eqref{eq:halt-def}.

\begin{theorem}[No total correct decider]\label{thm:no-decider}
iPRA proves $\forall d\,\neg\Dec(d)$.
\end{theorem}
\begin{proof}
Fix $d$ and suppose $\Dec(d)$. Its totality, applied to the well-formed
expression $C(e)$, supplies a finite terminating computation and a bit.
We eliminate this existential and split the decidable two-element case.

\emph{Output $1$.} By correctness, $\Id(C(e))$. By
\eqref{eq:calibration}, $\neg\Halt(e)$. But the supplied run and
\eqref{eq:diag-forward} give $\Halt(e)$, a contradiction.

\emph{Output $0$.} First prove $\neg\Halt(e)$. If a halting trace for
$e$ existed, \eqref{eq:diag-backward} would give a run of $d$ on $C(e)$
with output $1$. Determinism contradicts the supplied run with output
$0$. This proves $\neg\Halt(e)$ by implication to a contradiction.
Now \eqref{eq:calibration} gives $\Id(C(e))$. Correctness of the
observed output $0$ gives $\neg\Id(C(e))$, again a contradiction.

Both possible outputs contradict the assumptions, so $\neg\Dec(d)$.
Universal introduction gives the theorem.
\end{proof}

Only one instance of totality was used, namely at $C(e)$. In particular,
even a program that is required to terminate and be correct only on all
compiled inputs $C(z)$ is impossible. This proves the restricted-language
claim and hence the original MathWorld statement by inclusion of syntax.

\subsection{Constructive strength of the conclusion}

The proof did not reason from ``the diagonal program cannot fail to halt''
to a halting witness. In the first branch the witness is constructed from
an observed finite run. In the second branch an assumed witness is used
only to obtain a contradiction. Likewise, no inference
$\neg\forall q\,B(q)\Rightarrow\exists q\,\neg B(q)$ appears.

The result concerns every ordinary recursive algorithm, not just primitive
recursive candidate deciders. Primitive recursiveness is imposed on the
\emph{reductions, verifiers, and proof bounds}; the candidate $d$ is an
arbitrary machine code. Its assumed totality may have no primitive
recursive running-time bound, and none is used.

Finally, the theorem is not merely a failure to find a decision procedure,
and it is not a statement of independence of one particular expression.
It is an internally proved nonexistence theorem for a uniform algorithm.
There is no consistency premise in the displayed object-language theorem.
